\chapter{Skew et valeur relative de la structure par terme}\label{ch:s2:skew-and-term-structure-relative-value}

Entre septembre 2009 et septembre 2026, l'indice de volatilité à trois mois de Cboe se situait en moyenne 1.93~points au-dessus du VIX à un mois. Sur 7.6~\% des jours, la
hiérarchie s'est inversée, et le 12 mars 2020 l'indice à un mois se situait 18.23~points au-dessus de celui à trois mois. Deux options sur le même indice, une échéance à un mois
et une autre à trois mois, peuvent diverger sur un même avenir. Les traders de volatilité en valeur relative misent sur ce désaccord, ou sur la forme du smile,
plutôt que sur le niveau. Ce chapitre construit une surface synthétique dont le skew et la pente portent des erreurs de valorisation plantées, qui s'estompent, et les exploite avec
des options neutres en vega, couvertes en delta. Les erreurs de valorisation sont réelles~: les parties du P\&L qui en proviennent ont des ratios de Sharpe de 1.0 et 1.5. Les trades
rapportent 0.75 et 0.64, et l'essentiel de ce qu'elles rapportent vient d'ailleurs. La construction est \texttt{firm.volrv}.

\section{Trader le smile}

La volatilité implicite d'un indice est plus élevée pour les strikes bas que pour les strikes hauts~: c'est le skew (livre~1, chapitre~25). Cboe publie une mesure de celui-ci, le SKEW
est défini par $100 - 10S$, où $S$ est le prix de l'asymétrie du S\&P 500 calculé à partir d'un portefeuille d'options~; il monte quand la
queue gauche de la distribution risque-neutre s'alourdit. De janvier 1990 à septembre 2026, il valait en moyenne 123.2, avec des quantiles à 5~\% et 95~\% de 110.0 et 147.3.
Son plus bas, 101.23, est atteint le 19 mars 1991 et son plus haut, 183.12, le 18 février 2025. Le livre blanc de Cboe a relevé la faible corrélation entre les variations du SKEW et du VIX. Sur l'historique depuis
1990, les variations quotidiennes des deux sont corrélées à $-0.13$, et leurs niveaux à $-0.17$.

Bakshi, Kapadia et Madan ont comparé l'indice à ses composantes, à partir d'options sur le S\&P 100 et sur 30 de ses actions. Les distributions risque-neutres
des actions individuelles étaient bien moins négativement asymétriques que celle de l'indice. Ils attribuent l'asymétrie à l'aversion au risque et décomposent celle de chaque action en une part systématique
et une part idiosyncratique.

\begin{definition}[Trade de skew]\label{def:s2:skew-and-term-structure-relative-value:skew}
Un \emph{trade de skew}\index{trade de skew} achète des options à un strike et en vend à un autre, sur le même sous-jacent et la même échéance, dimensionné de façon que la
position n'ait aucun vega net et soit couverte en delta, afin de tirer profit d'une variation de l'écart entre leurs volatilités implicites plutôt que du
niveau.
\end{definition}

Le véhicule usuel est le risk reversal (livre~2, chapitre~19)~: long d'un call hors de la monnaie, short d'un put hors de la monnaie, ou l'inverse. Un trader
qui juge que l'aile de put est trop chère vend le put et achète le call. Le trade recèle une seconde exposition cachée. Une surface évolue avec le spot autant
qu'avec le temps~: sous une règle de sticky strike (livre~5, chapitre~7), chaque strike conserve sa volatilité quand le spot bouge, et sous une règle de sticky delta, le smile se déplace avec
le spot. Quand l'indice baisse, le put vendu se rapproche de la monnaie et son vega augmente. La position qui n'avait pas de vega en a maintenant, et elle est short.

\subsection*{La surface synthétique}

\texttt{firm.volrv} pose une surface quotidienne sur l'indice synthétique du chapitre~1. Les volatilités justes à un mois et à trois mois, à la monnaie, découlent de la
dynamique de la variance, prime incluse. Le coefficient de skew juste, c'est-à-dire la pente du smile par unité de moneyness standardisée $x = \ln(K/S)/(\sigma\sqrt\tau)$,
vaut $-0.12$ au niveau moyen. Il se raidit à mesure que le niveau monte~: 0.4 unité de skew pour chaque unité de volatilité à un mois. Par-dessus la surface juste, le
marché ajoute deux bruits plantés~: l'un sur le coefficient de skew (écart-type 0.06), l'autre sur la volatilité à trois mois (1.5 points de volatilité). Chacun est un
processus autorégressif de persistance quotidienne 0.98 et de demi-vie 34 jours. Ils représentent une demande qui pousse une partie de la surface puis s'estompe. Le
smile est $\sigma(x) = \sigma_{\mathrm{atm}}(1 + \text{skew}\cdot x + 0.02x^2)$.

Un trader ne voit pas les bruits. \texttt{features} extrait ce qui peut être observé~: le niveau (la volatilité à un mois), la pente (trois mois moins un mois) et
le skew. Il régresse ensuite la pente et le skew sur le niveau, en n'utilisant que les données antérieures à chaque jour, et conserve les résidus. Le skew brut est corrélé avec le
niveau à $-0.63$ et la pente brute à $-0.83$. Les ajustements sur l'échantillon complet donnent $-0.46$ pour le bêta du skew sur le niveau (la valeur plantée est $-0.4$) et $-0.20$ pour
celui de la pente. Les résidus restituent bien les bruits plantés~: ils sont corrélés à ceux-ci à 0.97 pour le skew et 0.92 pour la pente.

\section{Trader la structure par terme}

\begin{definition}[Trade de structure par terme]\label{def:s2:skew-and-term-structure-relative-value:term}
Un \emph{trade de structure par terme}\index{trade de structure par terme} achète de la volatilité à une maturité et en vend à une autre sur le même sous-jacent, dimensionné de façon que
la position n'ait aucun vega net et soit couverte en delta, afin de tirer profit d'une variation de la pente de la structure par terme plutôt que du niveau.
\end{definition}

Le calendar spread est la forme standard~: long à trois mois, short à un mois, sur des options à la monnaie de vegas égaux. Sa forme découle de la manière dont la
volatilité opère un retour à la moyenne. Quand la volatilité est élevée, le marché anticipe sa baisse, si bien que la volatilité à courte échéance se situe au-dessus de la volatilité à longue échéance et que la structure par terme s'inverse. Quand la volatilité est
faible, l'ordre se renverse. Une pente est donc, en grande partie, un niveau déguisé. Sur les indices réels (la \cref{fig:s2:skew-and-term-structure-relative-value:slope}),
la pente VIX3M moins VIX est corrélée au VIX à $-0.54$. Elle a valu en moyenne $-7.98$ points jusqu'en mars 2020. Elle revient aussi vite~: sa persistance
quotidienne est de 0.90, soit une demi-vie de 6.5 jours de bourse.

\begin{omfigure}
\begin{tikzpicture}
\begin{axis}[width=11cm,height=5.4cm,xlabel={\small année},ylabel={\small VIX3M $-$ VIX (en points)},tick label style={font=\footnotesize},
  xmin=2009.5,xmax=2027,ymin=-9,ymax=7,grid=major,grid style={gray!25},table/col sep=comma,/pgf/number format/1000 sep={}]
\addplot[omThm,thick] table[x=year,y=slope]{figdata/strategies-2/03-skew-and-term-structure-relative-value/slope.csv};
\addplot[black,thin,domain=2009.5:2027] {0};
\end{axis}
\end{tikzpicture}

\omcaption{La pente de la structure par terme du VIX~: l'indice de volatilité à trois mois de Cboe (VIX3M) moins le VIX, moyennes mensuelles, de septembre 2009 à septembre 2026.
Sous zéro, la courbe est inversée~: la volatilité à courte échéance se situe au-dessus de la volatilité à longue échéance. Série calculée à partir des historiques des indices de Cboe~; les séries brutes ne sont pas redistribuées.
Données~: \texttt{s2\_fetch\_term}.}\label{fig:s2:skew-and-term-structure-relative-value:slope}
\end{omfigure}

Johnson a étudié ce que dit la pente. La forme de la structure par terme du VIX renseigne sur le prix du risque de variance, et non sur les variations attendues
du VIX, ce qui rejette l'hypothèse des anticipations. Sa deuxième composante principale, une pente, prédit les rendements en excès des swaps de variance synthétiques sur le S\&P 500
des contrats à terme sur le VIX et des straddles sur le S\&P 500, à toutes les échéances. Un trader peut lire cela de deux façons. C'est un signal de timing pour les trades de
niveau du chapitre~1~: vendre de la variance quand la pente indique que la prime est élevée. C'est aussi une mise en garde pour un calendar spread, qui trade la pente elle-même.

\section{Entre sous-jacents liés}

\begin{definition}[Spread de volatilité entre sous-jacents]\label{def:s2:skew-and-term-structure-relative-value:cross}
Un \emph{spread de volatilité entre sous-jacents}\index{spread de volatilité entre sous-jacents} achète de la volatilité implicite sur un sous-jacent et en vend sur un
sous-jacent lié, comme deux indices, un indice et un fonds indiciel coté (ETF) sur celui-ci, ou un indice et ses composantes, en pariant que le ratio ou la différence des
deux volatilités revienne à sa plage habituelle.
\end{definition}

Cboe publie des indices de volatilité sur le Russell 2000 (RVX) et sur le Nasdaq-100 (VXN). Depuis septembre 2009, le RVX se situe en moyenne 5.2~points au-dessus du
VIX, avec un écart-type de 2.3 points et une demi-vie de 15.7 jours de bourse. Le VXN se situe en moyenne 3.0~points au-dessus, avec un écart-type de 2.3 et
une demi-vie de 15.3 jours. Les spreads reviennent plus lentement que la pente du VIX, et rien en eux n'est le niveau d'un seul marché déguisé. Cela en fait
le signal de valeur relative le plus net de ce chapitre. Il porte aussi le risque le moins familier. Les deux sous-jacents peuvent se séparer pour une raison (une
chute sectorielle, un changement dans la composition d'un indice), et le spread ne revient alors pas.

Personne ne trade les indices de volatilité eux-mêmes~: le trade porte sur des options sur chaque sous-jacent. Il comporte la couverture en delta de chaque sous-jacent, la fourchette de chaque
marché d'options, et la probabilité qu'une surface soit cotée moins souvent que l'autre. Le résultat de Bakshi, Kapadia et Madan est à la base du
trade entre sous-jacents le plus courant, skew d'indice contre skew d'action. Son risque est celui de la corrélation, abordé au chapitre~2.

\section{Les books et leur attribution}

Chaque mois, le book prend un trade unitaire dimensionné par l'opposé du signal, un z-score expansif écrêté à 2~: il vend ce qui est cher. Le trade est un risk
reversal d'une unité standardisée de part et d'autre de la monnaie, ou un calendar à un mois contre trois mois. Chaque jambe est dimensionnée de sorte qu'une hausse de 1~\% de
sa volatilité vaille une unité (0.18 point de volatilité pour une volatilité de 18). Chaque jambe est couverte en delta chaque jour et conservée jusqu'à la fin du mois. Ouvrir une jambe coûte
1~\% de sa volatilité, hypothèse du chapitre et non une fourchette sourcée. Le P\&L est attribué chaque jour par réévaluation dans l'ordre
(le \cref{lst:s2:skew-and-term-structure-relative-value:pnl})~: le spot et le temps (delta, gamma, theta, sur la surface du jour), puis le niveau (la surface juste
qui se déplace vers celle du lendemain), puis le bruit de skew, puis le bruit de pente. L'ordre compte~: le niveau est déplacé en premier, et les compartiments de skew et de pente ne
contiennent que les bruits plantés. Quatre books tournent sur 227 mois après une année d'initialisation~: chaque trade, signalé par la variable brute ou par son résidu sur le
niveau.

\begin{center}
\small
\begin{tabular}{@{}lccccccc@{}}
\toprule
& & & & \multicolumn{4}{c}{part du P\&L (\%)}\\
\cmidrule(l){5-8}
227 mois & Sharpe & asymétrie & pire (éc.-type) & spot, temps & niveau & skew ou pente & coût\\
\midrule
skew, brut & 1.13 & 1.91 & $-4.6$ & 94 & 5 & 9 & $-8$\\
skew, résidu & 0.75 & 1.48 & $-4.5$ & 73 & 24 & 16 & $-13$\\
calendar, brut & 0.07 & 4.74 & $-1.7$ & $-31$ & 160 & 155 & $-186$\\
calendar, résidu & 0.64 & 1.43 & $-2.4$ & 97 & $-16$ & 51 & $-32$\\
\bottomrule
\end{tabular}
\end{center}

Les effets injectés sont bien là. Pris isolément, le compartiment skew du livre de skew résiduel affiche un ratio de Sharpe de 1.01 et le compartiment de pente du
calendar résiduel affiche 1.51. Les livres bruts en captent moins~: 0.95 pour le compartiment skew du livre de skew brut et 0.84 pour le compartiment de pente du calendar brut. Mais ces
compartiments sont petits à côté du reste. Le bruit de skew injecté explique 16~\% du P\&L du livre de skew résiduel. Le livre de skew dont le signal repose sur la caractéristique brute
affiche le \emph{plus haut} ratio de Sharpe des quatre, et seulement 9~\% de son P\&L provient du retour du skew à sa moyenne. Le reste tient au spot et au temps. Le skew brut est
pentu quand la volatilité est élevée, si bien que ce livre vend l'aile put quand la volatilité est élevée. Ce qu'il encaisse, c'est la prime du chapitre~1, avec la paperasse d'un risk reversal.

Le calendar brut est l'autre cas. Son signal de pente est à 83~\% de niveau, si bien que son compartiment de niveau (160~\% du P\&L) et son compartiment de pente (155~\%) sont tous deux
grands, et les coûts comblent l'écart. Le calendar résiduel garde le compartiment de pente et se débarrasse de l'essentiel du niveau. La moitié de son P\&L est l'effet injecté.
L'autre moitié est le spot et le temps, et une partie de celle-ci relève aussi de la mauvaise valorisation~: une option achetée à bon marché rapporte davantage par le gamma qu'elle ne coûte en theta pendant qu'elle est
détenue. Mois après mois, le compartiment spot et temps est corrélé à la taille de la mauvaise valorisation en portefeuille, à 0.25 pour le livre de skew et à 0.14 pour le calendar. Le
reste est du bruit de gamma, et c'est lui qui domine la variance d'un mois à l'autre (la \cref{fig:s2:skew-and-term-structure-relative-value:calendar}).

\begin{omfigure}
\begin{tikzpicture}
\begin{axis}[width=11cm,height=5.4cm,xlabel={\small année},ylabel={\small P\&L cumulé (unités d'écart-type)},tick label style={font=\footnotesize},
  xmin=1,xmax=20,ymin=-15,ymax=50,grid=major,grid style={gray!25},table/col sep=comma,
  legend style={font=\scriptsize,at={(0.5,-0.32)},anchor=north},legend columns=4]
\addplot[black,thick] table[x=year,y=total]{figdata/strategies-2/03-skew-and-term-structure-relative-value/calendar.csv};
\addplot[omDef,thick] table[x=year,y=spot_time]{figdata/strategies-2/03-skew-and-term-structure-relative-value/calendar.csv};
\addplot[omThm,thick] table[x=year,y=slope]{figdata/strategies-2/03-skew-and-term-structure-relative-value/calendar.csv};
\addplot[gray,thick,dashed] table[x=year,y=level]{figdata/strategies-2/03-skew-and-term-structure-relative-value/calendar.csv};
\legend{total,spot et temps,bruit de pente,niveau}
\end{axis}
\end{tikzpicture}

\omcaption{Le P\&L du livre calendar résiduel et trois de ses compartiments, cumulés en unités de l'écart-type mensuel du livre (les coûts, non représentés,
comblent l'écart). Le bruit de pente injecté rapporte régulièrement. Le compartiment spot et temps est plus ample et bien plus bruité. Données~:
\texttt{s2\_volrv.books}.}\label{fig:s2:skew-and-term-structure-relative-value:calendar}
\end{omfigure}

Une seule graine peut flatter un livre. Sur huit graines portant sur le marché et sur la surface, les ratios de Sharpe moyens valent 0.90 pour le livre de skew brut (entre 0.47 et
1.26), 0.39 pour le livre de skew résiduel (0.12 à 0.75), $-0.07$ pour le calendar brut ($-0.43$ à 0.29) et 0.38 pour le calendar résiduel ($-0.10$ à
0.75). La graine du chapitre est du côté chanceux pour les deux livres résiduels.

\section{Risque~: ce qu'un livre de volatilité en relative value est réellement short}

Un livre sans vega net ni delta net à l'origine n'est pas un livre sans risque. Détenu chaque mois sans signal, une unité de chaque trade révèle
ce qu'il porte. Long de l'aile call et short du put, le livre rapporte 12.0 unités par mois au titre du spot et du temps, et perd 15.9 au titre du niveau~: $-6.2$ par mois en
tout. Son put vendu gagne du vega quand l'indice baisse, c'est-à-dire quand la volatilité monte. Long du calendar, le livre rapporte 9.6 unités par mois au titre du spot et du temps, puisque la jambe courte
à un mois encaisse la prime de variance par le theta, et perd 3.2 au titre du niveau~: 4.7 par mois en tout.

Le pire mois des deux livres de skew a commencé le jour 1\,470 (année 5.8). Les deux livres étaient à leur limite~: deux unités longues sur l'aile call et courtes sur le put. Le
krach est survenu neuf jours après la fin du mois. Pendant le mois, l'indice a progressé de 4.5~\% tandis que la volatilité à un mois passait de 15.2~\% à 9.7~\%. Le call long
monta jusqu'à son strike, et son vega augmenta au moment même où la volatilité qu'il détenait s'effondrait. Le mois a coûté 289 unités~: 119 au titre du spot et du temps (le theta du call dans un
rallye calme) et 157 au titre du niveau, soit 4.5 écarts-types mensuels du livre. Le signal résiduel n'avait aucun moyen de le voir venir~: le bruit de skew était
pentu ce mois-là et disait de vendre l'aile put. Un livre de volatilité en relative value est short de tout ce que ses jambes ne compensent pas. Pour un risk reversal, c'est un mouvement conjoint
du spot et de la volatilité dans un sens ou dans l'autre~: un krach quand il est short de l'aile put, un rallye calme avec une volatilité en baisse quand il est long du call. Pour un calendar
c'est le gamma du mois proche, que la jambe à trois mois ne couvre pas.

Trois règles en découlent. Attribuer avant de croire un ratio de Sharpe~: un \omterm{def:s2:skew-and-term-structure-relative-value:skew}{trade de skew} dont le compartiment skew ne représente qu'un dixième de son P\&L est un autre trade. Fonder le signal sur
les résidus par rapport au niveau, sinon le livre n'est qu'un trade de niveau avec des coûts en plus. Soumettre chaque livre à des chocs conjoints sur le spot, la volatilité et le skew, dans les deux sens, et non sur
l'un d'eux à la fois.

\section{Fiches de stratégie}

\begin{strategyfile}{Flattener de skew}\label{strat:s2:skew-and-term-structure-relative-value:skew}
\sfield{Qui paie, et pourquoi} Les hedgers qui achètent des puts sur indice au-delà de ce que justifient les autres points de la surface~; le paiement est la part de la cherté de l'aile put
qui s'estompe.
\sfield{Instruments et plateformes de négociation} Options cotées sur indice, puts et calls hors de la monnaie d'une même échéance.
\sfield{Signal} Le résidu du skew après sa dépendance au niveau, comparé à son historique.
\sfield{Dimensionnement et exécution} Risk reversals neutres en vega, couverts en delta chaque jour~; dimensionnés selon la perte du mois de krach, et non selon le ratio de Sharpe.
\sfield{Coûts} Deux fourchettes d'options à chaque roll et la couverture~; l'aile put est la moins liquide.
\sfield{Comment elle échoue} Un krach pendant qu'il est short de l'aile put~; un skew qui se raidit avec le niveau et ne revient jamais.
\sfield{Horizon, capacité, infrastructure} De quelques semaines à quelques mois~; un calibrage quotidien de la surface et un moteur de couverture.
\sfield{Backtester honnêtement} Attribuer le P\&L~; tester les signaux bruts et résiduels~; inclure les mois de krach.
\sfield{Sources} Bakshi, Kapadia et Madan (2003) sur le skew d'indice~; ce chapitre~: 0.75 pour une graine, 0.39 en moyenne sur huit, avec 16~\% du P\&L
provenant du skew.
\end{strategyfile}

\begin{strategyfile}{Spread de volatilité calendaire}\label{strat:s2:skew-and-term-structure-relative-value:calendar}
\sfield{Qui paie, et pourquoi} Une demande concentrée sur une seule échéance (les hedgers du mois proche, les émetteurs de produits structurés à plus longue échéance) qui s'estompe.
\sfield{Instruments et plateformes de négociation} Options cotées sur indice à deux échéances~; swaps de variance ou contrats à terme sur le VIX à deux maturités.
\sfield{Signal} Le résidu de la pente après sa dépendance au niveau.
\sfield{Dimensionnement et exécution} Vegas égaux~; couverts en delta~; jambe proche roulée chaque mois.
\sfield{Coûts} Deux fourchettes par roll~; la couverture de la jambe proche.
\sfield{Comment elle échoue} Un pic de volatilité, quand le gamma de la jambe proche vendue se retourne contre le livre~; un signal de pente qui n'est en réalité que le niveau.
\sfield{Horizon, capacité, infrastructure} De quelques semaines à quelques mois~; un calibrage de la structure par terme.
\sfield{Backtester honnêtement} Fonder le signal sur les résidus~; attribuer au niveau et à la pente~; inclure les inversions.
\sfield{Sources} Johnson (2017) sur la pente du VIX et les primes de risque de variance~; ce chapitre~: 0.64 pour une graine, 0.38 en moyenne.
\end{strategyfile}

\begin{strategyfile}{Spread de volatilité indice contre ETF}\label{strat:s2:skew-and-term-structure-relative-value:etf}
\sfield{Qui paie, et pourquoi} Des flux qui pèsent sur l'un de deux marchés d'options portant sur un sous-jacent quasi identique, comme les options sur indice face aux options sur un
fonds indiciel coté (ETF) qui réplique l'indice.
\sfield{Instruments et plateformes de négociation} Options sur indice contre options sur le fonds~; le suivi du fonds et ses dividendes fixent le spread juste.
\sfield{Signal} L'écart de volatilité implicite par rapport à l'écart juste ajusté du suivi.
\sfield{Dimensionnement et exécution} Neutre en vega, couvert en delta dans chaque sous-jacent.
\sfield{Coûts} Les fourchettes des deux marchés d'options~; des styles de règlement-livraison et d'exercice différents.
\sfield{Comment elle échoue} Des différences justifiées (style d'exercice, règlement-livraison, dividendes, impôts) prises pour une mauvaise valorisation.
\sfield{Horizon, capacité, infrastructure} De quelques jours à quelques semaines~; les deux surfaces en temps réel.
\sfield{Backtester honnêtement} Modéliser les différences de règlement-livraison et d'exercice avant de prendre un écart pour un signal.
\sfield{Sources} Aucun chiffre de performance vérifié.
\end{strategyfile}

\begin{strategyfile}{Spread de volatilité entre marchés}\label{strat:s2:skew-and-term-structure-relative-value:cross}
\sfield{Qui paie, et pourquoi} Une demande de protection sur un marché qui écarte sa volatilité de celle d'un marché apparenté.
\sfield{Instruments et plateformes de négociation} Options ou contrats à terme sur volatilité portant sur deux indices (petites capitalisations contre grandes capitalisations, une région contre une autre).
\sfield{Signal} Le spread ou le ratio de volatilité par rapport à son historique~; les indices RVX et VXN de Cboe face au VIX présentent un retour à la moyenne avec des demi-vies d'environ 15 jours.
\sfield{Dimensionnement et exécution} Jambes neutres en vega ou pondérées par le bêta.
\sfield{Coûts} Les fourchettes des deux marchés et les couvertures.
\sfield{Comment elle échoue} Un changement de régime sur un marché~: le spread se déplace et y reste.
\sfield{Horizon, capacité, infrastructure} Quelques semaines~; deux surfaces.
\sfield{Backtester honnêtement} Tester si le spread revient à la moyenne après un choc ou s'il suit une tendance.
\sfield{Sources} Historiques du RVX, du VXN et du VIX de Cboe (statistiques de ce chapitre)~; aucun chiffre de performance vérifié.
\end{strategyfile}

\section{Tutoriel~: même contrat à terme, deux prix}\label{tut:s2:skew-and-term-structure-relative-value:tut}

\begin{tutorial}
\textbf{Objectif.} Construire la surface synthétique, en extraire les caractéristiques, faire tourner les quatre livres et attribuer leur P\&L~; lire la structure par terme réelle, le SKEW et
spreads entre indices. \textbf{État final~:} le tableau et les deux figures.
\begin{tutsteps}
\item \textbf{Le trade et son attribution}.
\omcode{code/firm/volrv/firm_volrv.py}{104}{137}{Un trade unitaire conclu chaque mois, couvert en delta chaque jour, attribué par revalorisation séquentielle.}\label{lst:s2:skew-and-term-structure-relative-value:pnl}
\item \textbf{Signaux et livres}.
\omcode{code/strategies-2/03-skew-and-term-structure-relative-value/python/s2_volrv.py}{35}{64}{Z-scores expansifs et les quatre livres.}\label{lst:s2:skew-and-term-structure-relative-value:books}
\item \textbf{Exécuter} \texttt{s2\_fetch\_term.py} une fois, puis \texttt{table()}, \texttt{bucket\_sr()}, \texttt{seeds()}, \texttt{static()} et
\texttt{fig\_volrv.py}.
\end{tutsteps}
\textbf{Pour aller plus loin.} Couvrir l'exposition au krach du risk reversal avec un petit put acheté à plus longue échéance et en mesurer le coût~; changer la
règle de la surface pour passer à sticky delta et comparer les attributions~; ajouter une troisième maturité et trader la courbure de la structure par terme.
\end{tutorial}

\section{Construction~: volatilité en valeur relative}\label{bld:s2:skew-and-term-structure-relative-value:volrv}

\begin{build}
\bfield{Rôle} Une surface quotidienne avec des erreurs de prix de skew et de pente introduites volontairement~; les caractéristiques de la surface et leurs résidus sur le niveau~; des
risk reversals et calendars neutres en vega, avec attribution du P\&L.
\bfield{Interface} \texttt{SurfaceConfig(\dots)}, \texttt{surface\_path(r, v, cfg, scfg)}, \texttt{features(s, warmup)},
\texttt{trade\_pnl(r, s, kind, pos, cfg, scfg)}, \texttt{attribution(b)}.
\bfield{Règles} Résidus issus uniquement de régressions à fenêtre croissante~; jambes dimensionnées par le vega relatif à l'ouverture~; attribution dans un ordre fixe (spot et temps,
niveau, skew, pente) de sorte que les compartiments somment au total.
\bfield{Tests d'acceptation} \texttt{code/firm/volrv/tests/}~: un effet de niveau linéaire est supprimé exactement~; une surface immobile ne produit ni P\&L de niveau, ni de skew, ni de pente,
et les compartiments somment au total~; un risk reversal gagne lorsque le skew s'aplatit et perd lorsqu'il s'inverse.
\bfield{Pour aller plus loin} Dynamique sticky delta~; trades de courbure~; deux sous-jacents.
\end{build}

\begin{omsources}
\item G.~Bakshi, N.~Kapadia et D.~Madan, ``Stock return characteristics, skew laws, and the differential pricing of individual equity options'',
\emph{Review of Financial Studies} 16(1), 2003.
\item T.~L.~Johnson, ``Risk premia and the VIX term structure'', \emph{Journal of Financial and Quantitative Analysis} 52(6), 2017.
\item Cboe, SKEW Index white paper, 2011.
\item Cboe Global Markets, historiques quotidiens de VIX, VIX3M, SKEW, RVX et VXN.
\end{omsources}

\section{Exercices}

\begin{exercise}[$\star$]\label{exo:s2:skew-and-term-structure-relative-value:1}
Le prix de l'asymétrie du S\&P 500 vaut $S = -2.3$. Quelle est la valeur de l'indice SKEW~?
\end{exercise}

\begin{exercise}[$\star$]\label{exo:s2:skew-and-term-structure-relative-value:2}
Une jambe est dimensionnée de sorte qu'une hausse de un pour cent de sa volatilité vaille une unité. Sa volatilité vaut 20~\% et tombe à 19.6~\%. Que rapporte une position vendeuse sur la jambe~?
\end{exercise}

\begin{exercise}[$\star$]\label{exo:s2:skew-and-term-structure-relative-value:3}
Une autorégression a une persistance quotidienne de 0.98. Quelle est sa demi-vie en jours de négociation~?
\end{exercise}

\begin{exercise}[$\star\star$]\label{exo:s2:skew-and-term-structure-relative-value:4}
Pourquoi un calendar spread signalé sur la pente brute est-il largement un trade sur le niveau de la volatilité~?
\end{exercise}

\begin{exercise}[$\star\star$]\label{exo:s2:skew-and-term-structure-relative-value:5}
Johnson a montré que la pente du VIX prédit les rendements des swaps de variance, des contrats à terme sur le VIX et des straddles. Comment utiliseriez-vous ce résultat dans un trade sur le niveau, et que
signale-t-il à un trader de calendar~?
\end{exercise}

\begin{exercise}[$\star\star$]\label{exo:s2:skew-and-term-structure-relative-value:6}
Les distributions risque-neutre des actions individuelles sont moins fortement asymétriques négativement que celle de l'indice. Concevoir un trade sur cet écart et nommer son risque principal.
\end{exercise}

\begin{exercise}[$\star\star\star$]\label{exo:s2:skew-and-term-structure-relative-value:7}
\emph{Programmation.} Exécuter \texttt{table(0, 0.0)}, qui désactive les deux bruits introduits volontairement. Que se passe-t-il pour chaque livre, et qu'indique-t-il sur le
livre de skew brut~?
\end{exercise}

\begin{exercise}[$\star\star\star$]\label{exo:s2:skew-and-term-structure-relative-value:8}
\emph{Trouver la faille.} «~Notre livre de skew affiche un ratio de Sharpe de 1.13 et le skew revient à la moyenne, donc nous sommes rémunérés pour le retour à la moyenne du skew.~»
\end{exercise}

\section{Problème~: même future, deux prix}

\begin{problem}[{Problème du week-end --- valeur relative sur une surface}]\label{pb:s2:skew-and-term-structure-relative-value:1}
La surface synthétique du chapitre, les indices de Cboe et les données publiques.

\medskip\noindent\textbf{Partie I --- La surface.}
\begin{enumerate}
\item Définir un \omterm{def:s2:skew-and-term-structure-relative-value:skew}{trade de skew} et un \omterm{def:s2:skew-and-term-structure-relative-value:term}{trade de structure par terme}.
\item Comment l'indice SKEW de Cboe est-il défini, et quelles étaient ses statistiques~?
\item Qu'ont établi Bakshi, Kapadia et Madan sur les skews d'indice et d'actions~?
\item Comment le skew juste de la surface synthétique dépend-il du niveau, et quels sont ses bruits introduits volontairement~?
\end{enumerate}

\medskip\noindent\textbf{Partie II --- La structure par terme réelle.}
\begin{enumerate}[resume]
\item Donner la moyenne de la pente du VIX, la fréquence de son inversion et son extrême.
\item Pourquoi la pente est-elle corrélée au niveau~?
\item Qu'a établi Johnson~?
\item Donner les spreads du RVX et du VXN par rapport au VIX et leurs demi-vies.
\end{enumerate}

\medskip\noindent\textbf{Partie III --- Les livres.}
\begin{enumerate}[resume]
\item Comment les caractéristiques et leurs résidus sont-ils extraits~?
\item Comment chaque jambe est-elle dimensionnée, et dans quelles unités le P\&L est-il exprimé~?
\item Comment le P\&L est-il attribué, et pourquoi l'ordre compte-t-il~?
\item Donner les ratios de Sharpe des quatre livres et leurs attributions.
\end{enumerate}

\medskip\noindent\textbf{Partie IV --- Le verdict.}
\begin{enumerate}[resume]
\item Énoncer le \emph{résultat nommé}~: les ratios de Sharpe des \omterm{def:s2:skew-and-term-structure-relative-value:skew}{trades de skew} et de calendar, et la part de leur P\&L qui relève en réalité d'autre chose.
\item Que rapportent les effets introduits volontairement pris isolément~?
\item Que montraient les huit graines~?
\item Quel carry chaque trade porte-t-il lorsqu'il est tenu sans signal~?
\item Que s'est-il passé pendant le pire mois des livres de skew~?
\item Que se passe-t-il lorsque les bruits introduits volontairement sont désactivés~?
\item Comment effectuer un backtest honnête d'un livre de skew~?
\item En une phrase~: que vend à découvert un livre de volatilité en valeur relative~?
\end{enumerate}
\end{problem}

\section{Questions d'entretien}

\begin{interviewq}[$\star$ \iqroles{trader}]\label{iq:s2:skew-and-term-structure-relative-value:1}
Qu'est-ce qu'un risk reversal, et que vous dit son prix~?
\end{interviewq}

\begin{interviewq}[$\star\star$ \iqroles{trader}]\label{iq:s2:skew-and-term-structure-relative-value:2}
Vous êtes acheteur d'un calendar neutre en vega. Le marché chute brutalement. Que devient votre P\&L~?
\end{interviewq}

\begin{interviewq}[$\star\star$ \iqroles{chercheur}]\label{iq:s2:skew-and-term-structure-relative-value:3}
Comment sauriez-vous si un signal de skew est en réalité un signal de niveau déguisé~?
\end{interviewq}

\begin{interviewq}[$\star\star$ \iqroles{risque}]\label{iq:s2:skew-and-term-structure-relative-value:4}
Un livre de volatilité en valeur relative affiche un vega nul et un delta nul. Quels risques subsistent~?
\end{interviewq}

\begin{interviewq}[$\star\star$ \iqroles{développeur}]\label{iq:s2:skew-and-term-structure-relative-value:5}
Concevez une attribution quotidienne du P\&L pour un livre d'options. Que faut-il pour que les compartiments somment au total~?
\end{interviewq}

\begin{interviewq}[$\star\star\star$ \iqroles{chercheur}]\label{iq:s2:skew-and-term-structure-relative-value:6}
Dans le modèle de Black--Scholes, montrer que le P\&L d'une option couverte en delta sur un court intervalle vaut environ $\tfrac12\Gamma S^2(\sigma_r^2 -
\sigma_i^2)\,\mathrm{d}t$, et expliquer pourquoi deux jambes de même vega mais de volatilités implicites différentes n'ont pas un gamma égal.
\end{interviewq}
